% !TEX encoding = UTF-8
% !TEX TS-program = pdflatex
% !TEX root = ../tesi.tex

%**************************************************************
\chapter{Processi e metodologie}
\label{cap:processi-metodologie}
%**************************************************************

\intro{In questo capitolo vengono descritte le modalità con cui si è svolto lo stage e le tecnologie utilizzate.}\\

%**************************************************************
\section{Organizzazione del lavoro}

Il lavoro è stato svolto in parte a tempo parziale ed in parte a tempo pieno, data la concomitanza con il progetto di laboratorio del corso di Ingegneria del Sofware. 

Il piano di lavoro dello stage individuava tre periodi principali:

\begin{itemize}
    \item Un periodo di formazione sulle tecnologie utilizzate, della durata totale prevista di 80 ore e suddiviso in tre macro-categorie relative rispettivamente all'\acrshort{sso}, all'\acrshort{ssi} e a FreeIPA, Linux \acrshort{pam};
    \item Un periodo riservato alle integrazioni software, della durata totale prevista di 200 ore;
    \item Un periodo di stesura di documentazione relativa allo stage e alla tesi, della durata totale prevista di 20 ore.
\end{itemize}




\section{Tecnologie utilizzate}

Durante il periodo di stage ho avuto modo di utilizzare tecnologie per me nuove, oltre che approfondire la conoscenza di altre. Nelle sezioni seguenti sono illustrate tali tecnologie e strumenti.

\subsection{FreeIPA}
FreeIPA\cite{site:freeipa-website} (\autoref{fig:logo-freeipa}) è una soluzione open-source gratuita (GNU General Public License) di gestione dell'identità e dell'accesso per ambienti di rete basati su Linux/UNIX, originariamente sviluppato dalla comunità Fedora ed ora supportato da diverse organizzazioni, tra cui Red Hat e la FreeIPA Foundation. Consiste in un insieme di servizi integrati, che consentono di centralizzare l'autenticazione, l'autorizzazione e la gestione degli utenti e delle risorse in un'organizzazione.


\begin{figure}[!h] 
    \centering 
    \includegraphics[width=0.6\columnwidth]{logo-freeipa} 
    \caption{Logo di FreeIPA}
    \label{fig:logo-freeipa}
\end{figure}


FreeIPA è progettato per semplificare la gestione dell'identità e dell'accesso in ambienti di rete complessi, con molti utenti e computer. Consente agli amministratori di gestire facilmente l'accesso degli utenti a risorse e applicazioni, di delegare i privilegi di amministrazione e di definire ed applicare politiche di sicurezza coerenti in tutta la rete, come, ad esempio, limitare l'accesso alle risorse in base al ruolo dell'utente. Per fare ciò, mette a disposizione, oltre che agli strumenti della \acrshort{cli}, un'interfaccia utente web intuitiva per la gestione degli utenti, dei gruppi e delle risorse della rete. Inoltre, FreeIPA è altamente scalabile e può essere distribuito su più server per gestire grandi reti.

Per l'autenticazione degli utenti, FreeIPA utilizza il protocollo Kerberos: gli utenti possono accedere alle risorse della rete utilizzando le loro credenziali Kerberos, senza dover inserire le password ogni volta.
Per archiviare e gestire le informazioni sugli utenti, i gruppi e le risorse della rete, invece, utilizza il server di directory open-source 389 Directory Server, il quale offre funzionalità avanzate di ricerca, replica e sincronizzazione.

FreeIPA supporta l'autenticazione \acrshort{sso} tramite i protocolli \acrfull{saml} (Security Assertion Markup Language) e OICD: ciò significa che gli utenti possono accedere a più applicazioni utilizzando le stesse credenziali di accesso.

\subsection{Linux Containers}
\acrfull{lxc} è l'acronimo di Linux Containers, un sistema di virtualizzazione basato sul kernel Linux che consente di eseguire più sistemi operativi isolati su una singola macchina host. A differenza della virtualizzazione completa, in cui ogni sistema operativo guest ha accesso all'intero hardware dell'host, 

Utilizza la virtualizzazione basata sui contenitori, in cui ogni sistema operativo guest condivide le risorse hardware dell'host.
La condivisione del kernel fa sì che i container siano molto leggeri e veloci e che abbiano un overhead di risorse molto basso rispetto ad altre tecnologie di virtualizzazione.

\acrshort{lxc} fornisce un'interfaccia di riga di comando per la gestione dei container, che consente di creare, avviare, fermare, eliminare e gestire i container in modo semplice ed efficiente. Inoltre, supporta la creazione di immagini di container, che possono essere utilizzate per creare nuovi container in modo rapido e semplice.
\\
Durante l'attività di stage ho utilizzato principalmente container basati su immagini \acrshort{centos}, una distribuzione Linux basata su \acrshort{rhel} particolarmente adatta all'uso in ambiente server, che offre una vasta gamma di funzionalità e strumenti per gestire un'infrastruttura IT. 

In particolare, ho utilizzato l'ultima versione di \acrshort{centos} Stream (versione 9), che, a differenza della versione standard, riceve gli aggiornamenti in tempo reale durante lo sviluppo di \acrshort{ssi}, consentendo agli utenti di testare e fornire feedback sulle nuove funzionalità e correzioni di bug in anteprima\cite{site:centos-stream}.

Un'altra differenza tra \acrshort{centos} e \acrshort{centos} Stream è la durata del supporto: \acrshort{centos} ha, storicamente, fornito un supporto a lungo termine per le versioni rilasciate; la versione Stream, al contrario, è progettata per essere una piattaforma di sviluppo in continuo aggiornamento e non offre, dunque, il supporto a lungo termine.

\subsection{Secure Shell}
Secure Shell (\acrshort{ssh}) è un protocollo di rete crittografato utilizzato per la gestione sicura di dispositivi di rete e per l'accesso remoto a sistemi informatici. Il protocollo \acrshort{ssh} fornisce un canale di comunicazione sicuro tra due dispositivi, garantendo l'integrità, la riservatezza e l'autenticità delle informazioni trasmesse.

L'autenticazione avviene attraverso l'uso di chiavi pubbliche e private: in questo metodo, un'entità che desidera accedere a un sistema remoto genera una coppia di chiavi, una pubblica e una privata; la chiave pubblica viene fornita al sistema remoto, mentre la chiave privata viene conservata dall'entità; quando l'entità si connette al sistema remoto, la chiave privata viene utilizzata per autenticare l'entità.

Inoltre, \acrshort{ssh} utilizza la crittografia per proteggere i dati trasferiti tra i dispositivi. In particolare, il protocollo utilizza la crittografia a chiave simmetrica per proteggere i dati durante la trasmissione, e la crittografia a chiave pubblica per autenticare le parti coinvolte.

\section{Proxmox}

Proxmox è una piattaforma di virtualizzazione open source che combina la virtualizzazione basata su container (\acrshort{lxc}) e la virtualizzazione basata su macchine virtuali (\acrfull{kvm}) in un'unica soluzione. È progettato per consentire la creazione, la gestione e l'esecuzione di macchine virtuali e container su un unico sistema, fornendo un ambiente flessibile e scalabile per le infrastrutture IT.

Una delle caratteristiche principali di Proxmox è la sua interfaccia di gestione Web, chiamata \acrfull{pve}, che fornisce un'interfaccia grafica intuitiva per la gestione delle risorse di virtualizzazione. \acrshort{lxc} consente agli amministratori di creare, configurare e monitorare le macchine virtuali e i container, nonché di gestire l'archiviazione, la rete e altre risorse di sistema.

Proxmox offre, inoltre, funzionalità avanzate come la migrazione live delle macchine virtuali, che consente di spostare una macchina virtuale in esecuzione da un nodo all'altro senza interruzioni di servizio. Ciò offre flessibilità nella gestione delle risorse e consente di bilanciare il carico di lavoro tra i nodi di virtualizzazione.

La piattaforma è anche integrata con la gestione dell'archiviazione e della rete, consentendo la creazione e la gestione di storage e reti virtuali all'interno dell'ambiente di virtualizzazione: questo semplifica la gestione delle risorse e offre una maggiore flessibilità nella configurazione dell'infrastruttura di rete e storage.

% TODO: ridisegnare diagrammi